% !TEX root = ../main.tex
\section{Introduction}
\label{sec:intro}

Consider a decision you have to be able to defend two years from now.

A clinical operations team asks a language model system whether a compound is
approved for a paediatric indication. The system answers. Someone acts on the
answer. Later---in an audit, an inspection, a dispute---the question arrives:
what was that answer based on? Not whether the model is generally accurate. Which
document, in which version, as of which date, under which jurisdiction, supported
this specific claim.

Most current validation practice cannot answer that question, and the reason is
not that the practice is careless. It is that the practice measures the wrong
thing. The dominant approach evaluates a model by sampling questions with known
answers and scoring the outputs against them---\emph{reference-based evaluation},
in the terms of the natural-language-generation literature. It produces a
number---accuracy, or its complement, often reported as a hallucination
rate---and that number is then used as evidence of fitness for deployment.

The central argument of this paper can be stated in one sentence. \textbf{In
regulated settings, the relevant question is not merely whether an answer is
accurate. It is whether the organisation can reconstruct why this particular
assertion was produced, identify the authoritative and versioned evidence on which
it rests, establish that the evidence applies to the resolved case context, verify
the asserted relation, and show that reliance was authorised under the institution's
governing rules.} Accuracy is one input to that question and settles none of it.

This paper argues that reference-based evaluation cannot support the use it is put
to,
for two reasons that compound.

The first is a classification problem. Pooling every wrong answer into one
metric conflates failures with different causes and different remedies. A system
that faithfully reports an erroneous source has a data problem. A system that
fabricates a plausible-looking figure has a generation problem. A system that
correctly reports a regulation that was superseded last year, or that applies in
a different jurisdiction, has a context-resolution problem. These are three
separate defects, addressed by three separate teams, on three separate budgets.
One number does not tell you which one to spend.

The second reason is worse, and it is the argument of this paper.
Section~\ref{sec:canonical-form} shows that the system applies no canonical form
to its inputs. It computes a conditional distribution over token sequences given
token sequences; its inputs are strings; and two semantically identical
queries---a contraposed conditional, a quantity written in digits rather than
words, an active rather than passive phrasing---are two distinct points in input
space. Nothing in the architecture guarantees they map to the same output, and
measurement confirms they frequently do not. A database normalises before it
compares. A compiler canonicalises before it optimises. This system does
neither.

The consequence for validation is direct. If a test case passes, you have
learned that one string produced an acceptable output on one occasion. You have
not learned that the semantically identical string will, or that the same string
will next week, or that the same string will later in the same conversation. A
passing test does not generalise to the paraphrase, so the test suite does not
certify the behaviour it appears to certify. Accuracy is a property of a sample,
and this sample does not support the inference that validation requires of it.

That is a negative result, and on its own it would be unhelpful. The
constructive claim is that a different property does not have this defect.
\emph{Grounding}---whether an assertion is supported by evidence applicable to
the decision at hand, through a mechanism actually sensitive to that
evidence---is a property of the relation between an assertion and an evidential
base, not of a sampled output. It is stable under paraphrase in the way accuracy
is not, because it is a fact about the pipeline rather than about a draw from it.
It is therefore the property a governance regime can audit.

Section~\ref{sec:auditable} makes this precise, and proves the sharpest form of
the point: no accuracy figure identifies or bounds the grounding rate over the
unit square. No accuracy figure constrains the rate at which a system asserts
beyond its evidence, in either direction. A system can be highly accurate and
almost entirely ungrounded, and the second fact is the one that predicts what
happens when the questions change.

\paragraph{What this paper is not arguing.} Not that these systems are
unreliable in a way that precludes use. I build them, in exactly the regulated
settings this paper is about, and they work. Not that accuracy is uninteresting;
it is interesting, and it is the wrong thing to certify against. And not that
better models will not help. They will, on some axes and not others, and
Section~\ref{sec:architecture} says which.

\paragraph{Contributions.} Six, in the order they appear.

\begin{enumerate}
  \item \textbf{A redefinition of hallucination, indexed to a decision context.}
        Existing definitions ask whether an assertion is supported by a supplied
        source. This one asks the prior question---whether that source governs the
        case---by requiring support to be applicable at a resolved context of
        jurisdiction, valid time, material facts, and intended use, and by
        requiring the assertion to have been produced \emph{because of} the evidence
        rather than merely to agree with it. The consequence is a failure class that
        faithfulness and attribution metrics score as clean by construction, and a
        definition evaluable at the point of assertion without an answer key
        (Sections~\ref{sec:taxonomy}, \ref{sec:novelty}).
  \item A five-category taxonomy following from it---\textsc{Sound},
        \textsc{Lucky}, \textsc{Error}, \textsc{Fabrication}, and
        \textsc{Misapplied}---in which the fifth, assertion grounded in authentic
        but inapplicable evidence, is a distinct control failure with its own
        remedy and owner (Sections~\ref{sec:taxonomy}, \ref{sec:cell-v}).
  \item A proof that aggregate accuracy neither identifies nor bounds the grounding
        rate, so a validation programme reporting accuracy alone cannot distinguish
        radically different reliance profiles
        (Proposition~\ref{prop:independence}).
  \item An architectural account of why the failures are structural, grounded in the
        specific architecture and training regime these systems inherit, and an
        argument that the absence of a canonical form makes accuracy uncertifiable
        rather than merely incomplete
        (Sections~\ref{sec:architecture}, \ref{sec:canonical-form}).
  \item A set of nineteen invariants separating what assertion behaviour does not
        supply from what a system must be built to supply, each tied to a control
        (Section~\ref{sec:invariants}).
  \item A reference architecture for governed retrieval, a definition of
        institutional warrant, and an account of why the division of ownership
        across data, policy, and platform functions is itself a principal cause of
        failure (Sections~\ref{sec:governed-rag}--\ref{sec:three-teams}).
\end{enumerate}

\paragraph{Scope.} This paper is about one class of system: models that emit
natural-language assertions by autoregressive sampling over a token vocabulary,
fitted by next-token likelihood on large text corpora, and optionally post-trained
on human preferences. That is the decoder-only transformer language model
established at scale by the GPT-3 generation and everything built on that pattern
since. It is not about deep learning generally. Generative adversarial networks and
diffusion models are out of scope, as are encoder-only classifiers, rankers, and
extractors, and the conventional predictive models that make up most of an
enterprise analytics estate. Those components assert nothing, so grounding is not a
coherent question about them---which is not to say they are easy to validate;
Section~\ref{sec:not-easier} is explicit that they share several of the problems
described here. Section~\ref{sec:scope} states the boundary
precisely and says why it falls where it does.

\paragraph{Structure.} Section~\ref{sec:premise} states the premise the argument
rests on. Section~\ref{sec:architecture} identifies the systems in scope and
derives, from their architecture and training regime, why the failures this paper
describes are structural. Section~\ref{sec:taxonomy} redefines
hallucination as ungrounded assertion and gives the five categories that follow.
Section~\ref{sec:independence} proves that aggregate accuracy does not identify the
grounding rate.
Section~\ref{sec:grounding-context} defines the resolved context and the six
dimensions of grounding, and explains why retrieval and citation do not supply
them. Section~\ref{sec:canonical-form} develops the canonical-form argument and its
consequence for validation. Section~\ref{sec:invariants} states the invariants.
Section~\ref{sec:auditable} sets out why grounding rather than accuracy is
auditable, and where it can be measured. Sections~\ref{sec:governed-rag}
to~\ref{sec:three-teams} give the constructive architecture, the definition of
warrant with reliance tiers, and the organisational analysis.
Section~\ref{sec:evaluation} gives an evaluation protocol and
Section~\ref{sec:related} positions the work.

\paragraph{Notation.} Formal notation is used where it makes a claim checkable
and avoided where it would only make one look rigorous. Two operators, five
defined categories, one proposition with a two-line proof. Every definition is
stated over observable behaviour, so that any claim made here can be tested
against a running system.
